\chapter*{Ringkasan Proposal}
\addcontentsline{toc}{chapter}{RINGKASAN PROPOSAL}

% Ringkasan maksimum satu halaman, ditulis dengan jarak satu spasi.
\begin{singlespace}
  Mahasiswa kerap kesulitan mengelola tugas, jadwal kuliah, dan informasi
  akademik yang tersebar di berbagai kanal seperti email kampus, grup pesan,
  sistem pembelajaran daring, dan pengumuman dosen. Ketidakteraturan ini
  berdampak pada terlewatnya tenggat, keterlambatan pengumpulan tugas,
  meningkatnya prokrastinasi akademik, dan menurunnya kualitas belajar.
  Aplikasi manajemen tugas yang umum digunakan masih menuntut pengguna
  memasukkan seluruh data secara manual, padahal sumber informasi tugas
  terbesar mahasiswa justru berada di email kampus yang tidak terstruktur.
  Proposal ini mengusulkan pengembangan \textbf{Welya}, sebuah produk
  digital berupa asisten akademik berbasis kecerdasan buatan (AI) yang
  membantu mahasiswa mengelola tugas, jadwal, email kampus, dan kerja
  kelompok dalam satu tempat.

  Welya terintegrasi dengan Gmail, Google Tasks, dan Google Calendar melalui
  API resmi dengan otorisasi OAuth 2.0. Sistem memanfaatkan model bahasa
  besar (\emph{large language model}/LLM) untuk menganalisis email kampus
  secara otomatis, mendeteksi email yang mengandung penugasan, mengekstraksi
  atribut tugas seperti judul, mata kuliah, dan tenggat, lalu memecah tugas
  menjadi subtask beserta estimasi waktu pengerjaannya. Seluruh tenggat dan
  jadwal kuliah disajikan dalam satu dasbor dengan prioritas harian,
  dilengkapi kalender tersinkron, pencarian cerdas atas tugas, materi, dan
  catatan, serta ruang kerja kolaborasi untuk kerja kelompok. Kebaruan
  produk ini terletak pada deteksi tugas otomatis dari email kampus
  berbahasa Indonesia --- kemampuan yang belum tersedia pada aplikasi
  manajemen tugas populer seperti Google Tasks, Todoist, maupun Notion.

  Pengembangan dilakukan secara iteratif dengan metode \emph{prototype}
  melalui tahapan analisis kebutuhan, perancangan sistem, implementasi,
  pengujian, dan evaluasi kepada pengguna. Produk dibangun sebagai aplikasi
  web sekaligus aplikasi desktop sehingga dapat diakses lintas perangkat.
  Evaluasi mencakup pengujian fungsional \emph{black-box}, pengukuran
  ketepatan deteksi tugas dengan \emph{precision} dan \emph{recall}
  bertarget minimal 80\%, serta kuesioner \emph{System Usability Scale}
  (SUS) kepada minimal 10 responden mahasiswa dengan target skor di atas 68.

  Target luaran tugas akhir ini berupa aplikasi Welya (web dan desktop)
  yang berfungsi penuh dan siap diuji coba, laporan tugas akhir, artikel
  publikasi ilmiah, serta pendaftaran Hak Kekayaan Intelektual (HKI).
  Kehadiran Welya diharapkan membantu mahasiswa mengelola aktivitas
  akademiknya secara lebih teratur sekaligus menjadi contoh penerapan
  teknologi AI dalam mendukung produktivitas akademik di lingkungan
  perguruan tinggi.

  \vspace{0.5em}
  \noindent\textbf{Kata kunci}: asisten akademik, kecerdasan buatan,
  manajemen tugas, model bahasa besar, integrasi Google.
\end{singlespace}
\clearpage
